% ============================================================================
% Fichier  : annexes/Ann_B_outils.tex
% Titre    : Référentiel des Outils
% Version  : 1.0 -- novembre 2026
% ============================================================================
% Outils cités à titre d'exemples courants, sans recommandation commerciale ;
% versions et licences à vérifier avant usage. Tout outil de test technique
% ne s'emploie qu'avec une autorisation écrite (Chapitre 12, loi n° 2010/012,
% art. 69).
% ============================================================================

\chapter{Référentiel des Outils}
\label{ann:outils}

Les outils ci-dessous sont des exemples courants, cités pour situer les
tâches du cours. Aucun n'est imposé~; aucun ne remplace la méthode. Les
versions et les licences sont à vérifier avant usage.

\begin{boxalert}
\textbf{Règle.} Tout outil qui interroge un réseau ou un système (scan,
test d'accès, test d'intrusion) ne s'emploie qu'avec une autorisation écrite
du responsable du système, et de l'hébergeur s'il y a lieu
(Chapitre~\ref{chap:audit-technique}, loi n°~2010/012, article~69).
\end{boxalert}

\section*{Gestion et quantification des risques}

\begin{tableaucours}{L{3.6cm}L{5.0cm}c}{%
    \tableauHeader{Tâche} &
    \tableauHeader{Outils} &
    \tableauHeader{Chap.}
}
Registre, matrice, critères & Tableur partagé, à défaut d'un outil GRC &
\ref{chap:risque}, \ref{chap:grc-pratique} \\
Étude EBIOS Risk Manager & Modèle de tableur mis à disposition par
l'ANSSI~; outils labellisés par l'ANSSI & \ref{chap:ebios} \\
Simulation de Monte-Carlo & Python (bibliothèque numpy) ou tableur~; le
générateur \texttt{fair\_sanaga\_C05.py} de NDAP SUKULU & \ref{chap:fair} \\
Mesures $\mu$, $U$, $H$, $D_{\mathrm{KL}}$ & Tableur, Python &
\ref{chap:conclure}, \ref{chap:auditologie} \\
\end{tableaucours}

\section*{Audit et audit technique}

\begin{tableaucours}{L{3.6cm}L{5.0cm}c}{%
    \tableauHeader{Tâche} &
    \tableauHeader{Outils} &
    \tableauHeader{Chap.}
}
Sceller un export & Calcul d'empreinte SHA-256 (\texttt{sha256sum},
\texttt{Get-FileHash}) & \ref{chap:audit-technique} \\
Comparer à une référence & Guides CIS Benchmarks et outil d'évaluation
associé~; scripts de comparaison & \ref{chap:audit-technique} \\
Inventaire des services & Scanner de ports (par exemple Nmap) &
\ref{chap:audit-technique} \\
Vulnérabilités & Scanner authentifié (par exemple Greenbone OpenVAS)~;
calculateur CVSS du FIRST & \ref{chap:audit-technique} \\
Comptes et journaux & Extractions de l'annuaire et du collecteur de
journaux~; tableur pour les rapprochements & \ref{chap:audit},
\ref{chap:audit-technique} \\
Audit continu & Scripts de rapprochement planifiés, détenus par l'audit &
\ref{chap:emergents} \\
\end{tableaucours}

\section*{NDAP SUKULU}

Les fiches d'auto-évaluation et les ateliers du cours sont accessibles sur
NDAP SUKULU, sous le code \codeUE{}. Ils s'appuient sur des pièces et des
exports figés fictifs, générés pour le cours~: aucun test n'y porte sur un
système réel.
